\section{产品化指标与 ROI}

\subsection{指标清单}

部署 DSPy 系统后需要同步收集以下指标：LLM cost per query, tool invocation success rate, optimizer runtime, automation coverage（agent 占比）、user satisfaction 评分。

\begin{table}[H]
\centering
\caption{ROI 衡量指标}
\begin{tabular}{p{0.22\textwidth}p{0.28\textwidth}p{0.28\textwidth}}
\toprule
指标 & 观察窗口 & 期望方向 \\
\midrule
Cost / query & 每日 & 降低 10--20\% \\
Agent automation coverage & 每周 & 提升至 >70\% 工作量 \\
User satisfaction & 每次 release & 提升 0.2 分（满分 5） \\
Tool invocation success & 实时 & > 95\% \\
\bottomrule
\end{tabular}
\end{table}

\subsection{ROI 说明}

ROI 不只看节省的 API 费用，还要考虑 developer time reduction。通过 Evidence Matrix 的 trace 分析，可以量化 `optimizer improvement` 给业务带来的响应速率提升，从而算出 `value delivered / cost spent`。

\subsection{本章小结}

产品化指标确保 DSPy 的投入对齐业务目标，ROI 的计算需要把技术改善转化为 agent coverage 与用户满意度。

